\documentclass[11pt,a4paper]{report}

\usepackage{float} %dzieki temu pakietowi rysunki pojawiaja sie tam gdzie chce po uzyciu opcji H w srodowisku figure
\usepackage[font=small,labelfont=bf]{caption}
\usepackage{natbib}
\usepackage{import}
\usepackage[tight]{subfigure}
\usepackage{tikz}
\usepackage{pgfplots}
\usepackage{polski}
\usepackage[cp1250]{inputenc}
\usepackage{indentfirst}
\usepackage{color}
\usepackage{graphics}
\usepackage{geometry}
% ---- PAGE LAYOUT ----
\geometry{a4paper, top=1.5cm, bottom=1.5cm, left=1.5cm,
right=1.5cm, nohead}
\usepackage{amsmath,amssymb,amsfonts}
\begin{document}
\begin{center}
\large{\textbf{PRACA DOMOWA III}}
\end{center}

\noindent
imi� i nazwisko ........................................\\


\noindent
\textbf{Zadanie 1} Oblicz ca�ki nieoznaczone:\\
\textbf{(a)} $\int \frac{6x^3 + 1}{\sqrt{x^2 + 1}} dx$, 
\textbf{(b)} $\int \frac{x^2 - 3x + 2}{\sqrt{x^2 + 6x + 11}} dx$, 
\textbf{(c)} $\int\frac{2x^2 + 2}{\sqrt{-x^2 - 6x + 16}} dx$,  
\textbf{(d)} $\int\sqrt{x^2 - 6x + 12} dx$.\\

\vspace{0.25cm}

\noindent\textbf{odpowiedzi}:\\ 
\textbf{(a)} $\sqrt{x^2 + 1} (2x^2 - 4) + \ln|x + \sqrt{x^2 + 1}| + C$,  
\textbf{(b)} $\left(\frac x2 - \frac{15}{2}\right) \sqrt{x^2 + 6x^2 + 11} + 19\ln|x+3 +  \sqrt{x^2 + 6x^2 + 11}| + C$,\\ 
\textbf{(c)} $(x-9)\sqrt{-x^2 - 6x + 16} + 45\arcsin\left(\frac{x + 3}{5}\right) + C$ 
\textbf{(d)} $\left(\frac x2 - \frac 32\right) \sqrt{x^2 - 6x + 12} + \frac 32 \ln | x - 3 + \sqrt{x^2 - 6x + 12}| +C$.\\

\begin{figure}[H]
	\centering
	\begin{tikzpicture}
		\foreach \x in {0,...,34}
		\draw (0.5*\x,0)--(0.5*\x,20cm);
		\foreach \y in {0,...,40}
		\draw (0,0.5*\y)--(17cm,0.5*\y);
	\end{tikzpicture}
\end{figure}

\begin{figure}[H]
	\centering
	\begin{tikzpicture}
		\foreach \x in {0,...,34}
		\draw (0.5*\x,0)--(0.5*\x,27cm);
		\foreach \y in {0,...,54}
		\draw (0,0.5*\y)--(17cm,0.5*\y);
	\end{tikzpicture}
\end{figure}

\pagebreak


\noindent
\textbf{Zadanie 2} Oblicz ca�ki oznaczone:\\ 
\textbf{(a)} 1) $\int\limits_{1}^{4}\left( x^2 - \frac{1}{\sqrt{x}}\right)dx$, 
2) $\int\limits_{\pi/18}^{\pi/12}\sin(3x) dx$, 
\textbf{(b)} 1) $\int\limits_{2}^{3} \left(\frac{1}{x^2} - \frac{1}{x^3}\right)dx$, 
2) $\int\limits_{1}^{e^2} \frac{\ln xdx}{x}$,\\
\textbf{(c)} 1) $\int\limits_{0}^{1} \left(\frac{1}{x^2 + 1} - \frac{2}{1 - x^2} \right)dx$, 
2) $\int\limits_{1}^{2} \frac{x dx}{\sqrt{x^2 + 1}}$, 
\textbf{(d)} 1) $\int\limits_{1}^{4} \left(\frac{3}{x^2} \frac{2}{\sqrt{x}}\right)dx$, 
2) $\int\limits_{1}^{2} \frac{x dx}{(x^2 + 1)^2}$.\\


\vspace{0.25cm}
\noindent
\textbf{odpowiedzi}: \\
\textbf{(a)} 1) $19$,
2) $\frac{\sqrt{3} - \sqrt{2}}{6}$,  
\textbf{(b)} 1) $\frac{7}{72}$, 
2) $\frac{9}{2}$,  
\textbf{(c)} 1) $\frac{5}{4}\pi$,
2) $\sqrt{5} - \sqrt{2}$,  
\textbf{(d)} 1) $\frac{25}{4}$, 2) $\frac{3}{20}$.\\

\begin{figure}[H]
	\centering
	\begin{tikzpicture}
		\foreach \x in {0,...,34}
		\draw (0.5*\x,0)--(0.5*\x,21cm);
		\foreach \y in {0,...,42}
		\draw (0,0.5*\y)--(17cm,0.5*\y);
	\end{tikzpicture}
\end{figure}

\begin{figure}[H]
	\centering
	\begin{tikzpicture}
		\foreach \x in {0,...,34}
		\draw (0.5*\x,0)--(0.5*\x,27cm);
		\foreach \y in {0,...,54}
		\draw (0,0.5*\y)--(17cm,0.5*\y);
	\end{tikzpicture}
\end{figure}

\pagebreak

\noindent
\textbf{Zadanie 3}\\ 
\textbf{(a)} Wyznacz pole zawarte pomi�dzy liniami $y = 3x - x^2$ oraz $y = 2x$.\\ 
\textbf{(b)} Wyznacz pole zawarte pomi�dzy liniami $y = x^2 - 4x + 5$ oraz $y = 5 - x^2$.\\ 
\textbf{(c)} Wyznacz obj�to�� obszaru powsta�ego przez obr�t linii $y = 1-x^2$ wok� osi Ox, gdy $-1\leq x\leq 1 $.\\
\textbf{(d)} Wyznacz obj�to�� obszaru powsta�ego przez obr�t linii $y = \frac{1}{x^2}$ wok� osi Ox, gdy $1\leq x\leq 2$.\\


\noindent
\textbf{odpowiedzi}:\\ 
\textbf{(a)} $S = \frac 16$, \textbf{(b)} $S = \frac 83$, \textbf{(c)} $V = \frac{16}{15}\pi$,  \textbf{(d)} $V = \frac{7}{24}\pi$.\\

\begin{figure}[H]
	\centering
	\begin{tikzpicture}
		\foreach \x in {0,...,34}
		\draw (0.5*\x,0)--(0.5*\x,18cm);
		\foreach \y in {0,...,36}
		\draw (0,0.5*\y)--(17cm,0.5*\y);
	\end{tikzpicture}
\end{figure}

\begin{figure}[H]
	\centering
	\begin{tikzpicture}
		\foreach \x in {0,...,34}
		\draw (0.5*\x,0)--(0.5*\x,27cm);
		\foreach \y in {0,...,54}
		\draw (0,0.5*\y)--(17cm,0.5*\y);
	\end{tikzpicture}
\end{figure}

\end{document}